\section{Problem Formulation}
\label{sec:problem}

\subsection{Structured Harmful Meme Interpretation}
\label{sec:problem:task}

A meme is a multimodal instance
$\mathcal{M}=(I,T)$, where $I$ denotes the meme image and
$T$ denotes the text extracted from the image using optical
character recognition.
Conventional harmful meme detection predicts a binary label from
$\mathcal{M}$.
In contrast, we formulate meme understanding as a structured
interpretation problem that jointly identifies whether the meme is
harmful, who or what it targets, what communicative intent it
expresses, and which rhetorical and multimodal strategies construct
its meaning.

For each meme, the model predicts

\begin{equation}
\hat{\mathcal{Y}}
=
\left(
\hat{y}^{h},
\hat{y}^{p},
\hat{y}^{g},
\hat{y}^{i},
\hat{\mathbf{y}}^{r},
\hat{y}^{m}
\right),
\label{eq:structured_output}
\end{equation}

where $\hat{y}^{h}$ is the harmfulness prediction,
$\hat{y}^{p}$ denotes target presence,
$\hat{y}^{g}$ denotes target granularity,
$\hat{y}^{i}$ is the primary communicative intent,
$\hat{\mathbf{y}}^{r}$ is a multi-label prediction over rhetorical
tactics, and $\hat{y}^{m}$ denotes the multimodal relation through
which the image and text jointly convey meaning.

Harmfulness, target presence, target granularity, primary intent,
and multimodal relation are modeled as single-label classification
tasks.
Rhetorical tactic prediction is modeled as a multi-label task
because a meme may simultaneously employ multiple strategies, such
as sarcasm, stereotyping, dehumanization, or exaggeration.
Task-specific labels denoting insufficient or ambiguous supervision
are excluded using field-level masks, whereas a valid
\texttt{none} category is retained when the corresponding
phenomenon is absent.

\subsection{Evidence-Grounded Interpretation}
\label{sec:problem:evidence}

A structured label alone does not explain why a meme is considered
harmful.
We therefore require the framework to retain the evidence and
knowledge used to produce its predictions.
Let

\begin{equation}
\mathcal{E}^{\mathrm{int}}
=
\left\{
e^{\mathrm{text}},
e^{\mathrm{visual}},
e^{\mathrm{cross}}
\right\}
\label{eq:internal_evidence}
\end{equation}

denote internal evidence extracted directly from the meme.
It may include relevant OCR spans, visual entities or regions, and
cross-modal cues arising from the interaction between the image and
text.

The framework may additionally acquire a set of external contextual
candidates,

\begin{equation}
\mathcal{C}^{\mathrm{ext}}
=
\left\{
c_1,\ldots,c_K
\right\},
\end{equation}

from retrieval, context generation, or fallback mechanisms.
These candidates are not assumed to be reliable evidence.
A verifier evaluates their relevance, validity, and task-specific
support, producing a verified subset

\begin{equation}
\mathcal{E}^{\mathrm{ver}}
=
\mathcal{V}
\left(
\mathcal{M},
\mathcal{E}^{\mathrm{int}},
\mathcal{C}^{\mathrm{ext}}
\right)
\subseteq
\mathcal{C}^{\mathrm{ext}},
\label{eq:verified_evidence}
\end{equation}

where $\mathcal{V}$ denotes the verification procedure.
Retrieved candidates, generated hypotheses, fallback candidates,
verified evidence, and rejected candidates remain explicitly
distinguished through provenance metadata.

The final output therefore consists of both the structured
prediction in Eq.~\eqref{eq:structured_output} and an explanatory
record

\begin{equation}
\hat{\mathcal{A}}
=
\left(
\mathcal{E}^{\mathrm{int}},
\mathcal{E}^{\mathrm{ver}},
\mathcal{P},
\hat{R}
\right),
\label{eq:explanation_output}
\end{equation}

where $\mathcal{P}$ represents evidence and decision provenance and
$\hat{R}$ is a concise rationale grounded in the selected evidence.
The rationale is an explanatory output rather than a substitute for
the task predictions, and it should not introduce claims unsupported
by either the meme or the verified contextual evidence.

\subsection{Multi-Task Learning Objective}
\label{sec:problem:objective}

Let
$\mathcal{Q}=\{h,p,g,i,r,m\}$ denote the set of prediction tasks.
For sample $n$ and task $q$, let $m_n^q\in\{0,1\}$ be a supervision
mask indicating whether a valid label is available.
Let $w_n$ denote a sample weight derived from annotation-quality
metadata, and let $\lambda_q$ denote the task weight.
The overall learning objective is

\begin{equation}
\mathcal{L}(\theta)
=
\frac{1}{N}
\sum_{n=1}^{N}
w_n
\sum_{q\in\mathcal{Q}}
\lambda_q m_n^q
\,
\ell_q
\left(
f_\theta^q(\mathcal{M}_n),
y_n^q
\right),
\label{eq:multitask_objective}
\end{equation}

where $f_\theta^q$ is the prediction function for task $q$.
We use single-label classification losses for
$q\in\{h,p,g,i,m\}$ and a multi-label binary classification loss for
$q=r$.
This masked objective permits the framework to learn from partially
available structured supervision without treating missing or
ignored labels as negative instances.

\subsection{Cross-Dataset Generalization}
\label{sec:problem:generalization}

We study a source-to-target generalization setting.
The source collection

\begin{equation}
\mathcal{D}_{S}
=
\mathcal{D}_{\mathrm{COVID}}
\cup
\mathcal{D}_{\mathrm{Politics}}
\end{equation}

combines the COVID-19 and U.S.\ politics domains of HarMeme while
preserving their original domain identities.
It is partitioned into a source training set
$\mathcal{D}_{S}^{\mathrm{train}}$ and a source validation set
$\mathcal{D}_{S}^{\mathrm{val}}$.
The Facebook Hateful Memes collection forms a disjoint held-out
target set $\mathcal{D}_{T}$.

Model parameters are optimized only on the source training set:

\begin{equation}
\theta^{*}
=
\arg\min_{\theta}
\mathcal{L}
\left(
\theta;
\mathcal{D}_{S}^{\mathrm{train}}
\right).
\label{eq:source_training}
\end{equation}

All model-selection decisions, including early stopping,
hyperparameter selection, and decoding thresholds, are made using
$\mathcal{D}_{S}^{\mathrm{val}}$.
Let $\tau^{*}$ denote the resulting decoding policy.
Final performance is evaluated only after both $\theta^{*}$ and
$\tau^{*}$ have been fixed:

\begin{equation}
\hat{\mathcal{Y}}_{T}
=
f_{\theta^{*},\tau^{*}}
\left(
\mathcal{D}_{T}
\right).
\label{eq:heldout_evaluation}
\end{equation}

No target-domain sample is used for training, validation, threshold
selection, prompt construction, few-shot demonstration selection,
or retrieval-index construction.
This protocol evaluates whether harmfulness detection, structured
interpretation, and evidence-grounded reasoning learned from
HarMeme transfer to an unseen meme distribution.